% Drawing 83/10-1 -- folder 83, batch 10, archive page 198 (PDF page 199).
% Transcription: à droite, un parallélépipède en perspective : sommets marqués par des cercles, certaines arêtes renforcées à l'encre brune, d'autres en pointillé, flèches de rotation autour de deux axes ; étiquettes « $t\ldots$ », « $c$ ». À côté : « $\Sigma_a, \Pi_a$ ». Au-dessus, un petit losange et une flèche circulaire.
% Kind: parallelepiped in perspective with rotation arrows. Method 2 (schematic redraw from the facsimile). Status: partial.
% Notes: Parallelepiped: two vertices ringed (top and bottom), the left edge and two front edges heavy, hidden edges dotted; two curved rotation arrows on the faces; "t..." at the top right corner. Above: a small rhombus and a circular arrow. The labels c and Sigma_a, Pi_a are not in the crop (Sigma_a, Pi_a stands lower right); the brown ink of the transcription cannot be told on a grey render.
% Coordinates: 1 unit = 1 mm = 1/100 of the crop width (bbox in page fractions: [0.6, 0.06, 0.88, 0.26]).
% Redrawn by Claude (claude-opus-5-5), issue #73, 2026-11-07; not checked by a human; not compiled (no TeX in the session).
% Rights: reproduces part of the work -- behind RENDER_TRANSCRIPTS, never published by itself.
\documentclass[tikz,border=3pt]{standalone}
\usepackage{amsmath,amssymb}
\usetikzlibrary{patterns,arrows.meta}
\tikzset{pencil/.style={black!45}, struck/.style={black!35}, mark/.style={circle,draw,inner sep=1.3pt}, pt/.style={circle,fill,inner sep=0.9pt}, opt/.style={circle,draw,fill=white,inner sep=1.1pt}}
\begin{document}
\begin{tikzpicture}[x=1mm,y=1mm,line cap=round,line join=round,font=\scriptsize]
\draw (37,-57) -- (57,-46);
\draw (57,-46) -- (83,-51);
\draw (36,-83) -- (59,-90);
\draw (59,-90) -- (82,-72);
\draw[very thick] (37,-57) -- (36,-83);
\draw[very thick] (83,-51) -- (60,-62);
\draw[very thick] (50,-68) -- (82,-72);
\draw[densely dotted] (37,-57) -- (60,-62);
\draw[densely dotted] (83,-51) -- (82,-72);
\draw[densely dotted] (36,-83) -- (50,-68);
\draw (55,-63) -- (55,-70);
\draw (57,-46) circle[radius=1.5];
\draw (59,-90) circle[radius=1.5];
\draw[-{Stealth[length=2mm]}] plot[smooth, tension=0.55] coordinates {(44,-66) (46,-61) (50,-62)};
\draw[-{Stealth[length=2mm]}] plot[smooth, tension=0.55] coordinates {(60,-54) (63,-50) (66,-53)};
\node[inner sep=1pt] at (89,-49) {\tiny $t\ldots$};
\draw (21,-12) -- (27,-17) -- (21,-23) -- (15,-17) -- cycle;
\draw (16,-35) circle[radius=5];
\end{tikzpicture}
\end{document}
